\section*{Applications in additive combinatorics (original Section 3, contextual consequence)}
Let $\rho_k(N)$ be the largest density of a subset of $[N]$ without an arithmetic progression of length $k$. A corner is a triple of points $(x,y),(x+\delta,y),(x,y+\delta)$ for some $\delta\ne0$. Denote the largest density of a subset of $[N]\times[N]$ without a corner by $\rho_3^\angle(N)$.
As mentioned in the introduction, the NOF communication complexity of $f$ is closely related to corners theorems. Our Theorem 1 immediately implies,
\begin{theorem}
\[\rho_3^\angle(N)\ge2^{-2\sqrt{\log e}\sqrt{\log N}+o(\sqrt{\log N})}=2^{-2.4022\ldots\sqrt{\log N}+o(\sqrt{\log N})}.\]
There is an explicit corner-free subset of $[N]\times[N]$ of size
\[N^2/2^{2\sqrt{\log e}\sqrt{\log N}+o(\sqrt{\log N})}.\]
\end{theorem}
The derivation of Theorem 3 from Theorem 1 is an easy consequence of the following claim.
\begin{claim}[{[9], implicit}]
\begin{enumerate}
\item There is an optimal one-round protocol for the addition problem.
\item Let $T=T(x,y)$ be the message that the $P_z$ sends on inputs $(x,y)$ in a one-round protocol for the addition problem. Then the set
\[S(T)=\{(x,y):T(x,y)=T\}\]
is corner-free.
\end{enumerate}
\end{claim}
See [9,7,1,14,3] for more details about the above claim and the relation between communication complexity and additive combinatorics. The same comments and corollaries above apply also verbatim to the $(6,3)$ Theorem (e.g., [19]) and to the quantitative version of the triangle removal lemma.
